\subsection{Actual affine lines through rational points}
\label{subsec:actual-affine-lines}

\begin{definition}[Whole affine linear substitution]
\label{algebra:affine-line-substitution}
\lean{EnvelopeProperty.AffineLine.subst,
EnvelopeProperty.AffineLine.point,
EnvelopeProperty.AffineLine.aeval_comp_subst,
EnvelopeProperty.AffineLine.eval_subst,
EnvelopeProperty.AffineLine.eval_subst_zero,
EnvelopeProperty.AffineLine.subst_surjective,
EnvelopeProperty.AffineLine.aeval_zero_comp_subst}
\leanok
For a field $F$, any coordinate index type $\sigma$, and actual vectors
$a,b\in F^\sigma$, construct the genuine $F$-algebra homomorphism
\[
 \ell_{a,b}^*:F[U_i:i\in\sigma]\longrightarrow F[T],
 \qquad U_i\longmapsto a_i+b_iT.
\]
Evaluation at every parameter $t$ equals evaluation of the entire original
polynomial at $a+tb$. At zero this recovers the whole rational point;
for an arbitrary actual rational-point algebra homomorphism, its coordinate
values give exactly that original homomorphism. If $b\ne0$, the whole
polynomial homomorphism is genuinely surjective.
\end{definition}
\begin{proof}
Use the polynomial algebra universal property with the actual linear
coordinate polynomials. Algebra homomorphism extensionality proves the
entire evaluation-composition identity; this preserves coefficients as
well as every polynomial. For a nonzero direction choose an actual
coordinate $i$ with $b_i\ne0$. The original polynomial
$b_i^{-1}(U_i-a_i)$ maps exactly to $T$. All coefficients are also in the
image; actual polynomial induction therefore gives surjectivity on every
polynomial, over the given field. No mere coordinate-value
inference for seminorms is used.
\end{proof}

\begin{theorem}[An actual nonconstant line meets a principal open]
\label{algebra:affine-line-avoids-polynomial}
\lean{EnvelopeProperty.AffineLine.exists_eval_ne_zero,
EnvelopeProperty.AffineLine.exists_direction_subst_ne_zero,
EnvelopeProperty.AffineLine.exists_direction_parameter_ne_zero,
EnvelopeProperty.AffineLine.finite_bad_parameters}
\leanok
\uses{algebra:affine-line-substitution}
Assume $F$ is infinite and $\sigma$ is nonempty. Given any rational point
$a$ and any nonzero polynomial $h\in F[U_i:i\in\sigma]$, there is an
actual nonzero direction $b$ such that $\ell_{a,b}^*(h)\ne0$.
Only finitely many rational parameters of this line lie in $V(h)$.
In particular the constructed line through $a$ is not contained in the
closed hypersurface; it meets the genuine principal open $D(h)$.
\end{theorem}
\begin{proof}
Polynomial function extensionality over the infinite field produces an
actual rational point $z$ with $h(z)\ne0$. If $h(a)=0$, take $b=z-a$.
The unequal polynomial values force $b\ne0$, and evaluation at parameter
one makes the restricted polynomial nonzero. If $h(a)\ne0$, choose a
coordinate of the nonempty index set and its unit direction. Evaluation
at zero already proves the restricted polynomial nonzero. The roots of
an actual nonzero univariate polynomial form a finite set; the full
evaluation identity identifies these roots with the bad rational
parameters. Since the field is infinite, a parameter outside that finite
set gives the claimed meeting with $D(h)$.
\end{proof}

\begin{theorem}[The genuine closed line and its original scheme point]
\label{scheme:affine-line-closed-immersion}
\lean{EnvelopeProperty.AffineLine.schemeMap,
EnvelopeProperty.AffineLine.schemeMap_isClosedImmersion,
EnvelopeProperty.AffineLine.schemeMap_zero}
\leanok
\uses{algebra:affine-line-substitution,algebra:affine-line-avoids-polynomial}
The actual contravariant Spec map of this whole polynomial homomorphism
is a genuine closed immersion when $b\ne0$. The actual scheme point of
$\operatorname{Spec}F[T]$ at zero composes with this morphism to the
original whole rational point of the affine polynomial scheme.
For the preceding direction it meets $D(h)$.
\end{theorem}
\begin{proof}
Take the actual Spec map of the constructed polynomial ring homomorphism.
Its proved surjectivity supplies the genuine scheme closed immersion.
Contravariant Spec functoriality turns the full evaluation-at-zero
homomorphism equality into the stated whole scheme diagram. This is
actual scheme algebra, with no independently supplied line or incidence
certificate.

Over the genuine complete algebraically closed field constructed earlier,
infinitude is automatic, so this supplies the rational-point and
principal-open line ingredient for Fang--Han--Xia Proposition~5.22.
Passing to an open affine subset, its later curve dimension arguments,
and reduction, normalization or subharmonicity are distinct producers;
none is assumed or asserted by these line theorems.
\end{proof}
